\section{Summary}\label{sec:agg-creds-summary}

This chapter considered anonymous credentials in which each attribute is
certified by its own issuer, and issuers coordinate neither with each other nor
with the registration authority or verifiers. Users combine credentials from
independent issuers into a single aggregate credential and present it with one
proof.

The registration authority binds each user's key to a tag, and is not involved
after registration. Every anonymous credential scheme already assumes that
users' keys are certified, and we make this role explicit. The registration
authority learns only keys and tags, so the role adds no trust. Binding
credentials to the tag lets issuers stay uncoordinated, and since each
user has one tag, issuers need not remember whom they have signed for.

The instantiation combines the $\AGHO$ SPS scheme, our tag-based aggregate
signature scheme $\ATS$, and proofs based on Sigma protocols. With $n = 128$ issuers and $k =
64$ disclosed attributes, a presentation takes 20.1\,ms to produce and
31.1\,ms to verify.

Several directions remain open. Proving that every blinded hash uses the same
$\delta$ in $\group{T}$ would make presentations constant-size, at the cost of
one pairing per disclosed attribute for the user. The GGM is needed only for the
registration signature, so a rerandomisable SPS scheme under standard
assumptions would remove it, at the cost of larger registration signatures. A
concrete security analysis of the scheme would give its security at real
parameters.

\rebekah{kept Extensions in the construction; the summary only names revocation
as future work}
Epoch-based revocation, one of the extensions of
\Cref{sec:agg-creds-construction}, covers attributes with natural expiry dates.
Revoking a whole user through their tag, or a single attribute credential, needs
more than this, and we leave it to future work.

The scheme uses the issuer public keys in the clear. As the identity of the
issuers can reveal information about the user, recent work has focused on
issuer hiding. Adding issuer hiding to this scheme is also future work.
